\chapter{Reflection}\label{ch:reflection}

In programming, reflection refers to any case where the program needs to reason about itself, especially during execution.
This can be anything from the list of all the classes that implement a specific interface, or the list of fields inside one of those classes.
This can be extremely useful when serializing data into a human-readable format, as in these cases we would need to iterate all the data members inside the object.

Reflection is a common capability of most high level languages, for example, JS even has the ability to print out the code for a class of function.
In C\# and Java, we have the ability to list every class or my type of members inside them.

This ability is missing from C++ by default as it is Wii complex and would be extra data only useful in a few use cases.
Because of this everyone who needs this data in their code needs to find some way to generate and store the information.

If reflection is not available then it is possible to use manually written functions that return the same data.
<Code example>

In the example, we have a base class that has a virtual function that all derived classes must overwrite and return the class metadata.
This approach works and can give all the data that one could need.
However, this is a lot of manual work and puts the responsibility on the coder to write the function correctly.
This takes time away from actual development.

Another solution is to use a tool that can generate these functions or generate the data in a portable format that can be loaded during runtime.


\section{Why we need reflection}\label{sec:why-we-need-reflection}

In a game engine reflection is used for (among other things) serialising the level data.
For this we need to know what data fields an object has and witch of those need to be written to the level file.
On top of this, we also need to display all the fields in the editor, so the designer can change those with ease.

\subsection{Other engines}

The two engines that use the easiest to understand process for getting reflection data are UE and Source

\subsubsection{Source}
Source uses the simplest approach where the developer is responsible for creating a function that returns all the metadata about a class.
To make this easier they have made a few macros that complete in the necessary lines of code.
<Code example>

As source uses an external editor that doesn't depend on the game source code they also create a description file that tells the Hammer editor what entities exist in the game with what possible properties.

\subsubsection{UE}

Ue uses a tool that generates the reflection information by parsing the source code of the game and engine.
This tool is strongly integrated into the build process and relies on a few changes to the compiler.
Partial


\section{Shadow header tool}\label{sec:shadow-header-tool}

For Shadow Engine I have decided to try and create a similar tool to what UE has.
This would allow the users of the project a more seamless experience and free up development time in the future.
I have explored multiple possible ways to create such a tool, as for this I would need a partial AST (Abstract syntax tree).
This is notoriously hard to generate for Cpp.

I have found multiple online tutorials on how to create compiler plugins for clang.
Clang is a compiler frontend for the famous LLVM compiler framework, and is one of the 3 most used C/Cpp compilers.
While these tutorials showed great results and allowed full integration into the compilation process, they almost all required compiling clang from scratch.

I have tried following some of these tutorials, but I have run into many issues that would have need way more time to solve than I had available.
Most of my problems boiled down to one thing, there is no distribution of clang for Windows that supports loading and developing plugins.
Because of this I would need to compile clang myself and this is not an easy task and can differ based on the system being used.

After realising that getting clang plugins working is a way bigger task than I should work on, I have decided to write a simple parser in C\#.
I have chosen C\# because I'm familiar with it and has better string manipulation methods than cpp.

...

Cpp attributes

\subsection{Goal}
Get a list of all the classes that are in the code.
To be expandable in the future the tool parses all the classes no matter if they have an attribute that marks it for reflection.

\subsection{Ast}
After the initial regex based implementation I got inspiration and initial help creating a more advanced code processor that is much closerin design to a classical compiler. These software usually have the following steps:
- Lexer
- parser
- etc..
The lexer is responsible for dividing the source code into lexicographical tokens.
These are usually things like keywords, identifiers, and the different kinds of brackets.
This is done by iterating over the input text character by character and if the character is a known char we put it into a list of tokens.
<Example code>
In my implementation I parse the following tokens:
<List>
These are not all the possible tokens, but they are enough for what this tool needs to understand.

\subsection{Parser}
Parsing is the step where we iterate over all the tokens that makeup a file and form logical units out of them.
These are going to be the Nodes that make up the final AST. For each file the root of the corresponding AST is going to be a CompilationUnitNode.

For example the following series of tokens represent a basic class.

<Example of class tokens>

Would turn into this:

<Pseudo code of class node>

During parsing we find these patterns and form node from them.

To get this we can do a similar forward iteration as we did while lexing but this time we iterate on the token level.
Most of these patterns are well documented in the cop specifications and are more easily available online in cppreference.com

In the tool I have created a partial parser as I do not care about most of the cop code.
As a result my parser ignores all the function bodies and any declarations that it doesn't recognise as a class head or a member field.

\subsection{Output}
Both of the parsers (regex based and lexing) output the data they gather in a common format.
This makes it easy to add new code processors in the future, thus leaving the possibility of a clang integration open.

The output writer is a simple function that finds all classes with the correct attribute (SH::Reflect) and,
using simple string interpolation writes the necessary functions into a common file.
